\documentclass[11pt,a4paper]{article}
\usepackage[utf8]{inputenc}
\usepackage[T1]{fontenc}
\usepackage{amsmath,amssymb}
\usepackage{graphicx}
\usepackage{hyperref}
\usepackage[numbers]{natbib}
\usepackage{geometry}
\geometry{margin=1in}

\title{Daily Paper Summary: Macaron-V1 --- Towards Open Continual Learning\\with Self-Improvement and Mixture-of-LoRA}
\author{Internbot}
\date{August 12, 2026}

\begin{document}
\maketitle

\section{Paper Summary}

\subsection{Problem and Motivation}

Current large language models are trained on massive datasets but become static once deployed---they cannot learn from real-world interactions, adapt to individual users, or improve over time. Multi-task joint training of a single model creates destructive interference where improving one capability degrades another (termed ``gourd-like mutual interference''). Mind Lab~\cite{macaron2026} introduces Macaron-V1, a system architecture for \emph{open continual learning} that addresses both problems through Mixture-of-LoRA (MoL) and a recursive self-improvement (RSI) framework called MindForge. ``Open'' refers to the extensibility of the architecture: new specialist adapters can be registered without modifying or retraining the base model, enabling modular, collaborative capability growth.

\subsection{Method}

\paragraph{Mixture-of-LoRA (MoL).} Macaron-V1 freezes the entire base model (GLM-5.2 at 744B parameters) and trains small, independent LoRA adapters for each capability domain. Four specialist adapters are defined: L0 (Chat, routing entry point and conversational backbone), L1 (Agent, personal-life tasks and tool use), L2 (Coding, SWE and repository workflows), and L3 (GenUI, UI rendering and action generation). Each adapter is ${\sim}1$B parameters (${\sim}0.13\%$ of the base), yielding ${\sim}748$B total versus ${\sim}2{,}976$B for replicated full models---approximately \textbf{26\% of the replicated layout footprint}.

\paragraph{Three-stage routing lifecycle.} Each user turn follows a three-phase process: (1)~\textbf{Route}: L0 classifies the incoming request into one of four canonical labels using constrained decoding with a 24-token budget, implemented as a \texttt{change\_model} tool call; (2)~\textbf{Answer}: the selected specialist responds using its own ``conversation view''---full reasoning history for its own turns, but collapsed 192-token summaries for other specialists' turns; (3)~\textbf{Summary}: the specialist emits a 192-token summary shared as cross-turn context for other specialists. Routing accuracy is \textbf{99.12\%} on a 6{,}448-sample production trace. Per-hop latency is 4.68s on the Venti variant (0.54s route + 3.17s answer + 0.97s summary), with no detected quality degradation.

\paragraph{Catastrophic forgetting via architectural isolation.} Rather than algorithmic regularization (EWC, experience replay), Macaron-V1 prevents forgetting by construction: the frozen base preserves all pre-trained knowledge, and isolated LoRA adapters cannot interfere with each other since they occupy disjoint parameter spaces. Improving one specialty (e.g., coding) cannot degrade another (e.g., chat). The paper validates that text-only adapter training retains the base model's multimodal capabilities.

\paragraph{MindForge: recursive self-improvement.} A three-stage improvement cycle drives continual learning: (1)~\textbf{Discovery}: starting from seed tasks, the model constructs harder tasks, verifies answers, and filters for quality and learning value, identifying base-model failures; (2)~\textbf{Expansion}: the model solves selected tasks and iteratively searches context configurations (prompts, tools, skills, HCP files) to address failures---adaptive search reached \textbf{122/122} selected base-failure tasks versus only 11/122 for the strongest single-configuration sweep; (3)~\textbf{Update}: successful configurations and trajectories are distilled into adapter parameters through the MinT post-training platform. A key philosophical commitment is treating the \emph{harness} (runtime environment) as a first-class optimization target alongside model weights, standardized through a Harness Context Protocol (HCP) that ensures training-serving parity.

\paragraph{Training infrastructure.} LongStraw enables long-context RL by evaluating the shared prompt once into a reusable resident state and replaying only response branches that learn (2.1M-token GRPO on 8$\times$H20 GPUs). R3 (Rollout Routing Replay) records expert routing decisions at rollout time and masks tokens following misaligned paths from gradient computation. IcePop monitors token-level probability ratios and zero-weights tokens outside trusted bands. LoRA-RL training on the 744B model uses 64 H800 GPUs at ${\sim}10\%$ of full-parameter RL cost.

\subsection{Results}

On internal and public benchmarks, Macaron-V1 (Venti) leads on ChatBench (58.3 vs.\ GPT-5.5's 55.5), LivingBench (64.0 vs.\ 61.9), PinchBench (94.0 vs.\ 91.8), TerminalBench~2.1 (87.6 vs.\ 83.4), and UI4ABench (87.8 vs.\ 72.1---a 20.7-point improvement over the GLM-5.2 base). On coding benchmarks, it is competitive but trails frontier models: SWE-Verified 85.6 (vs.\ Opus~4.8's 88.6), DeepSWE 58.4 (vs.\ GPT-5.5's 70.0). Importantly, Macaron-V1 improves over its own base model on \emph{every} benchmark, validating that MoL adds value without degrading capabilities.

A collaboration experiment with 200 independently-trained LoRA adapters shows ${\sim}25\%$ improvement over a single-adapter approach, with performance following a log-linear relationship as adapter count increases. Cross-LoRA KV cache reuse experiments show no significant quality degradation (0.632$\pm$0.019 vs.\ 0.650$\pm$0.030 on VitaBench).

\subsection{Limitations}

The full recursive self-improvement cycle (where Update feeds back into Discovery for compounding gains) has not been evaluated across multiple generations---whether continual learning actually compounds remains an open question. The current evaluation isolates the Expansion stage with a frozen base and no optimizer steps. Several leading benchmarks (ChatBench, LivingBench, UI4ABench) are internally developed, raising evaluation objectivity concerns. The coding performance gap versus frontier models (DeepSWE $-11.6$ points behind GPT-5.5) suggests the architecture has not yet closed the gap for deep SWE tasks. All experiments use a single base model family (GLM-5.2/Qwen~3.6); generality across architectures is untested.

\subsection{Relevance to Research Topics}

Macaron-V1 represents a significant architectural stance on continual learning for LLM agents: rather than training a single model that must balance all capabilities, it proposes modular adapter-based isolation where each capability occupies its own parameter space. This directly contrasts with our recent review of the ``SFT Conflicts, RL Coexists'' paper~\cite{zhu2026sft}, which showed that RL training produces approximately orthogonal updates across tasks within \emph{shared} parameters. Macaron-V1 enforces orthogonality by construction through architectural separation, trading off potential cross-task synergies for guaranteed non-interference. The MindForge RSI framework connects to our prior reviews of self-evolving agent systems (SkillOS, SkillOpt, EvoArena), and the harness-as-optimization-target philosophy extends the model-scaffolding co-design ideas seen in DCAS and SEED.

\section{Follow-up Research Directions}

\subsection{Direction 1: Bridging Architectural and Gradient Orthogonality for Continual Multi-Task RL}

\paragraph{Gap.} Macaron-V1 enforces task isolation through separate LoRA adapters (architectural orthogonality), while ``SFT Conflicts, RL Coexists''~\cite{zhu2026sft} shows that GRPO naturally produces near-orthogonal updates within shared parameters (gradient orthogonality). These represent two extremes: full isolation (no cross-task transfer) versus full sharing (emergent orthogonality). Neither approach characterizes the optimal trade-off---isolated adapters cannot benefit from cross-task synergies, while shared parameters may lose orthogonality under distribution shift or longer training.

\paragraph{Approach.} Design a \emph{soft-isolation} continual RL framework with a shared LoRA trunk plus task-specific LoRA branches. The shared trunk captures cross-task synergies while the branches preserve task-specific knowledge. During GRPO training, monitor the cross-task gradient interference metric $I_{\text{RL}}(i,j) \leq V_i \cdot V_j$ from the SFT-RL coexistence theory; when interference exceeds a threshold $\tau$, automatically split the shared trunk into task-specific branches for the offending parameters. This creates an adaptive architecture that starts shared (maximizing transfer) and progressively isolates (preventing interference) as needed. The Geometry Conflict framework~\cite{geoconflict2026} could provide the per-parameter interference detection mechanism.

\paragraph{Feasibility.} Requires implementing interference-aware parameter splitting during RL training, which adds moderate overhead (one extra backward pass for interference estimation). The key risk is that the splitting criterion may be too conservative (isolating too early, losing transfer) or too permissive (isolating too late, allowing forgetting). Calibration against both Macaron-V1's full-isolation baseline and the SFT-RL coexistence shared-parameter baseline would provide clear reference points.

\subsection{Direction 2: MindForge-Style RSI for On-Policy Distillation Curricula}

\paragraph{Gap.} MindForge's Discovery stage generates training curricula by identifying base-model failures and constructing harder tasks. Current on-policy distillation (OPD) methods use static teacher-student configurations without adaptive curriculum generation. The student's training data is determined by the teacher's rollout distribution, which may not target the student's specific failure modes. Combining MindForge's failure-driven curriculum with OPD could create a self-improving distillation pipeline.

\paragraph{Approach.} Implement a three-stage OPD-RSI loop: (1)~\textbf{Discover}: evaluate the student on a capability benchmark to identify failure clusters; (2)~\textbf{Distill}: generate teacher rollouts specifically targeting failure clusters, using Flux-OPD's contextual correction~\cite{flux2026opd} to anchor distillation to the student's current capability; (3)~\textbf{Verify}: measure whether the student's failure rate on the targeted clusters decreases without degrading performance on previously mastered clusters (using the variance-limited interference bound from~\cite{zhu2026sft} as a stability criterion). The loop iterates, with each generation targeting a fresh set of failure clusters. CoRT's counterfactual replay~\cite{cort2026} could provide token-level credit assignment within the distillation loss to focus learning on the most informative tokens.

\paragraph{Feasibility.} Each component exists independently (failure detection, targeted OPD, interference monitoring); the contribution is their integration into a closed-loop curriculum. Main risk is that failure-targeted distillation may overfit to specific failure patterns without generalizing. The 122/122 coverage result from MindForge's Expansion stage suggests that adaptive targeting is highly effective, but this was achieved through harness optimization rather than weight updates---whether the same coverage translates to parameter-level distillation is untested.

\subsection{Direction 3: Adapter Collaboration Scaling Laws for Continual Agent Learning}

\paragraph{Gap.} Macaron-V1's 200-adapter collaboration experiment shows ${\sim}25\%$ improvement with log-linear scaling, but this result is preliminary. The fundamental questions remain unanswered: What is the theoretical scaling law for adapter collaboration? Does the log-linear relationship hold at 1{,}000+ adapters? How should adapters be \emph{selected} for composition when the catalog grows to millions (as MinT supports)? And critically, how does adapter collaboration interact with the continual learning setting---when new adapters are added over time, do older adapters need to be retrained?

\paragraph{Approach.} Formalize adapter collaboration as a submodular function maximization problem: given a catalog of $N$ adapters, select a subset $S \subseteq [N]$ of size $k$ that maximizes a composite benchmark score $f(S)$. If $f$ is submodular (diminishing returns), greedy selection achieves a $(1 - 1/e)$ approximation. Empirically validate submodularity by measuring $f(S \cup \{a\}) - f(S)$ for random subsets of increasing size. Then derive a scaling law $f(k) = \alpha \log k + \beta$ and test whether it holds across adapter counts from 10 to 1{,}000. For the continual learning dimension, measure whether adapters trained on earlier data distributions degrade when composed with newer adapters, and whether periodic re-routing (updating L0's classification boundaries) suffices to maintain composition quality without adapter retraining.

\paragraph{Feasibility.} Requires access to a large adapter catalog and systematic composition experiments, which are computationally expensive but methodologically straightforward. The submodularity test is cheap (requires only evaluation, not training). The main challenge is the combinatorial explosion of composition subsets at large $N$; random sampling with confidence bounds can mitigate this. Low theoretical risk, moderate computational cost.

\section{Conclusion}

Macaron-V1 presents a comprehensive system architecture for open continual learning that addresses catastrophic forgetting through architectural isolation (frozen base plus independent LoRA adapters) and post-deployment stagnation through recursive self-improvement (MindForge). The Mixture-of-LoRA design achieves 99\% routing accuracy with 26\% of the storage footprint of replicated models, while improving over the base model on every benchmark. The model-harness co-design philosophy---treating the runtime environment as a first-class optimization target---is a notable conceptual contribution that extends beyond the specific MoL architecture. While the full closed-loop RSI cycle remains unvalidated and coding performance trails frontier models, the paper establishes a principled framework for continual agent improvement. Combined with our recent review of gradient orthogonality in multi-task RL~\cite{zhu2026sft}, this opens a research direction on the optimal trade-off between architectural and emergent parameter-space isolation for continual multi-task learning.

\bibliographystyle{plainnat}
\begin{thebibliography}{10}

\bibitem{macaron2026}
Mind Lab (75 contributors).
\newblock Macaron-V1: Towards Open Continual Learning with Self-Improvement and Mixture-of-{LoRA}.
\newblock \emph{arXiv preprint arXiv:2608.09819}, 2026.

\bibitem{zhu2026sft}
Kejian Zhu, Zhuoran Jin, Shangqing Tu, Hongbang Yuan, Yushi Bai, Kang Liu, Juanzi Li, and Jun Zhao.
\newblock {SFT} Conflicts, {RL} Coexists: A Theoretical and Empirical Analysis of Multi-Task Learning for {LLMs}.
\newblock \emph{arXiv preprint arXiv:2608.03573}, 2026.

\bibitem{geoconflict2026}
Geometry Conflict Authors.
\newblock Explaining and Controlling Forgetting in {LLM} Continual Post-Training.
\newblock \emph{arXiv preprint arXiv:2605.09608}, 2026.

\bibitem{flux2026opd}
Flux-{OPD} Authors.
\newblock On-Policy Distillation with Evolving Contexts.
\newblock \emph{arXiv preprint arXiv:2607.28022}, 2026.

\bibitem{cort2026}
{CoRT} Authors.
\newblock Counterfactual Replay for Token-Level Rubric-Guided Policy Optimization.
\newblock \emph{arXiv preprint arXiv:2607.25659}, 2026.

\end{thebibliography}

\end{document}
